\section{Introduction}\label{sec:current-overview}\label{sec:introduction}
A predictive state records everything about a past experiment that can
affect its future laws. Its linear dimension need not measure the memory
required to realize those laws with nonnegative transition probabilities.
The distinction matters in causal experiments: commands can change what is
observed next, observation changes the state on which later commands act,
and the resources used to retain or communicate information must themselves
be counted. This paper studies that relation between predictive compression
and exact finite realization, within a broader theory of causal resource
comparison and optimal discrimination.

Two forms of observation will be used. A terminal response assigns a
probability row $b(w)$ to every finite command word $w$. An instrument
response assigns a joint probability to every finite command--outcome
history. A classical realization has one fixed collection of stochastic
tables and a finite set of live labels. Unless a mixed-start comparison is
explicitly stated, the initial label is pure and included in the count.
The tables may use arbitrary real probabilities; a lower bound therefore
cannot be evaded by increasing numerical precision. Their description and
execution costs are separate resources.

\subsection*{Reversible responses and the positivity gap}
Let $V_b$ be the span of residual functions $b^w(z)=b(wz)$, with dimension
$d$. Suppose that the command shifts are invertible on $V_b$ and generate
a group with compact closure $G_b$. This compact reversibility condition
holds for prepared finite-dimensional unitary command families after all
directions invisible to future tests have been removed. The finite
realization criterion, Theorem~\ref{thm:r72-reversible}, says that an exact
finite classical realization exists precisely when $G_b$ is finite.

Theorem~\ref{thm:r73-normal-form} gives the resource structure of every
such realization. It replaces continuing stochastic dynamics by a finite
permutation action and at most one transient initializer, without increasing
the number of live labels. The original transition matrices need not form
a group. The proof compresses their observable action through recurrent
classes of words which act as the identity on $V_b$. If $N$ is the number
of deterministic residuals and $\kappa_G(b)$ the minimum size of a
probability factorization $b(g)=\lambda P_gD$, then
\[
 m(b)=\min\{N,1+\kappa_G(b)\}.
\]
The bounded decoder and the initial mixture are essential. The least size
$\nu_G(V_b)$ of a permutation module containing $V_b$ is a lower bound
for $\kappa_G(b)$, and can be strictly smaller.

We determine how this positivity gap changes under terminal readout
contrast. Let $\pi$ be the group-averaged probability row and put
$b_t=\pi+t(b-\pi)$, $0\le t\le1$. For every finite group and every
finite outcome set, Theorem~\ref{thm:r74-low-contrast} proves
\[
 \kappa_G(b_t)=\nu_G(V_b),\qquad
 m(b_t)=\min\{N,1+\nu_G(V_b)\}
\]
for all sufficiently small positive $t$. Throughout $t>0$, the group,
module and residual cardinality stay unchanged. The proof constructs a
strictly positive preparation in a smallest containing permutation module
and controls the decoder explicitly. Theorem~\ref{thm:r74-contrast-thresholds}
then gives a complete finite filtration by attained thresholds. Its general
algebraic computability is distinct from a model-specific closed form.

For every group whose subgroups are normal, including all finite abelian
groups, Theorem~\ref{thm:r74-quotient-lp} makes the optimization linear.
Within-orbit preparation can be absorbed into a quotient decoder; only a
mixture among distinct quotient orbits remains. Theorem~\ref{thm:r74-quotient-dual}
gives necessary and sufficient payoff inequalities, strict exclusion
certificates and exact contrast endpoints. Rational data have rational
optimal controllers and rational thresholds. For abelian groups, character
support determines the small-contrast module bound. At higher contrast,
positivity adds a genuine obstruction: a four-command binary family has
$d=6$ and sixteen residuals at every $t>0$, but its minimum is eleven
labels for $t\le3/5$ and thirteen for $t>3/5$. Explicit payoff witnesses also
give a robust neighbourhood of the larger-memory regime. The retained
prepared-qubit $S_3$ model has an exact four-to-six transition at $2/3$.

\subsection*{Observed instruments and conditional memory}
Reversibility is only one way in which response memory is constrained.
The retained Theorem~\ref{thm:r68-instrument} treats known finite quantum instruments
on $\mathbb C^h$ whose every positive normalized branch state is at least
$\eta I/h$. For one fixed known prepared state it constructs a simulator
with at most
\[
 1+\left\lceil(4h/\eta)^{h-1}\right\rceil
\]
live labels. This theorem was proved before the present revision, together with sparse
branch rows, several-step interiority and a state-positive noise exponent.
A finite physical dictionary contains the whole posterior
region; branchwise barycentric expansions intertwine each completely
positive map with a nonnegative outcome matrix. Consequently the same
tables reproduce all adaptive visible histories, terminal POVMs, almost surely finite stopped
laws and nonnegative history costs. No state-independent outcome
probabilities or commuting-command assumption is required. Applying
depolarizing noise after each observed branch supplies the hypothesis for
arbitrary finite instruments. A complementary common-contraction theorem
gives an explicit finite positive realization of supplied linear responses
under one norm certificate for all commands.

Average mixing is insufficient for the instrument conclusion.
Theorem~\ref{thm:r74-depolarizing-obstruction} constructs a stationary
qubit family with an exactly depolarizing unconditioned channel, uniformly
positive outcome probabilities, one-dependent observations and Hankel rank
three. If its conditional phase has order $q\ge3$, its exact classical
memory is $q$ with a free mixed start and $q+1$ with a counted pure start.
A non-root-of-unity phase requires infinite memory. The upper uses one
invariant polygon for all outcomes; the lower uses an observable cylinder
whose relative peripheral spectrum forces a $q$-cycle. Equality in the
cycle bound is rigid enough to rule out a pure $q$-label initialization.
Thus neither linear rank, unconditioned mixing time nor finite dependence
determines the positive realization resource. The exact finite and infinite
cases lie in one fixed physical dimension and one fixed outcome alphabet.

\subsection*{Classical ingredients and contribution}
Finite residual equivalence, stochastic automata, positive realization and
observable operator models are classical
\cite{RabinScott59,Heller65,CarlylePaz71,Jaeger00,ThonJaeger15,MonrasWinter14}.
Schwarz, Flor and Brown developed the structure of nonnegative and
stochastic matrix groups \cite{Schwarz64,Flor69,Brown64,Schwarz62}.
Symmetric cone factorizations and extended formulations
\cite{GPT13,KPT10}, nonnegative-rank perturbations
\cite{BocciCarliniRapallo11}, and restricted nonnegative factorization
\cite{GillisGlineur12} precede the present work. Their general mechanisms
are credited. The assertions here concern the exact observable normal form
with its initial charge, universal attainment of the permutation-module
bound at small contrast, the normal-quotient linear reduction, and the
resulting exact resource transitions.

Invariant-polytope simulation of quantum instruments and adaptive
composition are explicitly developed by Okay~\cite{Okay26}; pure-state
nets and convex approximation of density matrices are established tools
\cite{AubrunSzarek17}. The earlier instrument results of this paper supply the posterior condition,
quantitative physical dictionary and exact resource accounts, with complete
proofs retained in Section~\ref{sec:interior-instruments}. They are used
here as established premises, not counted again as new results.
Finite-dimensional quantum processes without finite classical realizations
already occur in positive-realization theory \cite{Fanizza24}.
Holroyd and Liggett's finitely dependent colorings provide earlier
one-dependent non-HMM processes \cite{HolroydLiggett16}. We claim neither
first phenomenon. The specific assertion is the simultaneous rank-three,
uniformly positive, exactly depolarizing qubit construction and its sharp
$q$, $q+1$, or infinite classification. Classical finite-group recognition,
real-algebraic minimization and finite-word equivalence also remain credited
\cite{Serre07,BabaiBealsRockmore93,KieferWachter14,Tzeng92}.

\subsection*{Relation to the causal resource theory}
The main article retains the full actual-history comparison theorem and
common paid construction. Its lower witness uses the competitor's actual
pre-call histories. The paired supplement retains performed-law
compactness, variational and dual attainment, robust tied optimality,
normalized stopping rigidity, finite-width transfer, rank-constrained
discrimination, finite-order quantum measurements, two-stage reset control,
memory and communication accounts, spectral geometry and optimal strategy
implementation. It also contains the periodic partition, leakage and
arithmetic laws, all with their original hypotheses and complete proofs.
An independently compiled companion retains absolute calibration and
nonintervening directional sensing.

These layers have distinct quantifiers. General causal comparison applies
to its declared experiments; the group theorem classifies compactly
reversible responses; the posterior theorem supplies a sufficient
condition for general instrument families; the exact examples solve their
specified models; and local geometric results remain local. The theorem
and preservation maps make those dependencies explicit. None of these
results is asserted to complete independent A1, A2 or further Theta
construction objectives.
